\section{The Quantum}

The strangeness of quantum mechanics, the probabilities and the correlations that no classical picture can match, comes
out of the deterministic dynamics plus the symmetry of the states, with no randomness added by hand. This is one of the
strongest results in the program, because a deterministic base reproducing quantum behavior is far harder, and far more
telling, than one that quietly injects quantum randomness. The pieces are the Born rule, the Bell correlation, and the
nature of measurement.

The Born rule first, the rule that turns an amplitude into a probability by squaring it, the exponent two that everyone
memorizes and no one derives. Here it is forced. Ask for a law that assigns probabilities to alternatives and is
additive over them in the right way, consistent with the symmetry of a state entangled with its environment, and the
only consistent assignment makes the probability the square of the amplitude. The exponent comes out two, not put in. An
experiment confirms the square law from the additivity, the weight of an outcome scaling as the square of its amplitude,
with the exponent measured rather than assumed.

The Bell correlation next, the sharpest divide between the quantum and the classical. Two parts of an entangled state
can be correlated more strongly than any local classical theory allows, up to a precise ceiling, the Tsirelson bound,
two times the square root of two. On the mesh the dynamics that hops a tone is the spin-exchange interaction, and the
exchange interaction is an entangling gate. An experiment takes a plain product state of two charge-qubits, applies the
exchange dynamics, and produces a maximally entangled state whose correlations violate the classical bound and reach
exactly the Tsirelson ceiling, two root two, while a product-state control stays classical. So the quantum ceiling is
saturated by the substrate's own dynamics, not posited. A quantum walk on the mesh also spreads at constant speed where
a classical random walk only crawls as the square root of time, the ballistic signature of genuine quantum motion, which
an experiment measures directly.

So the Born rule is forced by additivity, the Tsirelson bound is reached by the exchange dynamics, and the quantum
spreading is ballistic, three quantum signatures coming out of a deterministic reversible rule. What remains is the
measurement problem, what happens when an outcome becomes definite, taken up next.
